\chapter{Cristals Kyber}
\section{Introdução}

%Da pra fazer um teste praticoooo: https://github.com/pq-crystals/kyber
% A nist usa, mas nào achei docs deles, vai ter que fazer só o teste pratico e coemntar quea  nist está adotando pra PQC, SE VIRAAAA!!! Comentario dee IA: e que o kyber é um dos candidatos a ser o padrao de criptografia pós-quântica do nist, e que é baseado em problemas de aprendizado com erros (LWE) e é considerado resistente a ataques de computadores quânticos.


O CRYSTALS-Kyber é um sistema de criptografia de chave pública baseado em problemas de aprendizado com erros (LWE) e é considerado resistente a ataques de computadores quânticos. 
Ele é parte do projeto CRYSTALS (Cryptographic Suite for Algebraic Lattices) e é um dos candidatos para o padrão de criptografia pós-quântica do NIST (National Institute of Standards and Technology). O CRYSTALS-Kyber é projetado para ser eficiente em termos de tempo de execução e tamanho de chave, tornando-o uma opção viável para aplicações práticas de criptografia pós-quântica.
\section{Estrutura do CRYSTALS-Kyber}
O CRYSTALS-Kyber é baseado em um esquema de criptografia de chave pública que utiliza matrizes e vetores para realizar operações de criptografia e descriptografia. Ele é composto por três algoritmos principais: geração de chaves, criptografia e descriptografia.